\chapter{Primary Operators and Finite Conformal Transformations}
\label{ch:volume-iii-primary-operators-finite-transformations}

\section{Spin Representations and Primary Data}

Let \(\rho\) be a finite-dimensional representation of \(SO(D)\) on a vector
space \(V_\rho\).  A local operator valued in this representation is written
\[
  \mathcal O_a(x),
  \qquad
  a=1,\ldots,\dim V_\rho .
\]
The infinitesimal spin matrices are denoted
\[
  (S_{\mu\nu})_a{}^b=-(S_{\nu\mu})_a{}^b .
\]
They obey the \(SO(D)\) Lie algebra in the representation \(\rho\), in the
convention in which a finite rotation \(R\) near the identity is represented as
\(\rho(R)=\exp(-\frac{\ii}{2}\omega^{\mu\nu}S_{\mu\nu})\).  Thus the
\(S_{\mu\nu}\) are Hermitian spin matrices in a unitary representation.

A conformal primary operator of spin \(\rho\) and scaling dimension
\(\Delta\) is an operator whose state at the origin obeys
\begin{equation}
  \widehat K_\mu\ket{\mathcal O_a}=0,\qquad
  \widehat D_{\rm rad}\ket{\mathcal O_a}=\Delta\ket{\mathcal O_a},\qquad
  \widehat J_{\mu\nu}\ket{\mathcal O_a}
  =
  -(S_{\mu\nu})_a{}^b\ket{\mathcal O_b}.
  \label{eq:primary-state-data}
\end{equation}
The translation generators create the descendant states:
\[
  \widehat P_{\mu_1}\cdots\widehat P_{\mu_k}\ket{\mathcal O_a}
  \longleftrightarrow
  \partial_{\mu_1}\cdots\partial_{\mu_k}\mathcal O_a(0).
\]
Thus, under the state--operator correspondence, a positive-energy conformal
module is generated from a lowest-energy state by repeated action of
\(\widehat P_\mu\).  Under the assumptions used in this volume, namely a
reflection-positive radial Hilbert space with discrete finite-multiplicity
eigenspaces of \(\widehat D_{\rm rad}\), the local operator space decomposes
as a direct sum of such primary modules after null states are quotiented.

An operator at a general point is obtained by translating the insertion from
the origin, equivalently by its Taylor expansion in descendants:
\begin{equation}
  \mathcal O_a(x)
  =
  \sum_{k=0}^\infty
  \frac1{k!}
  x^{\mu_1}\cdots x^{\mu_k}
  \partial_{\mu_1}\cdots\partial_{\mu_k}\mathcal O_a(0).
  \label{eq:operator-from-origin-translation}
\end{equation}
In the Hermitian translation convention
\([\widehat P_\mu,\mathcal O(x)]=\ii\partial_\mu\mathcal O(x)\), the
corresponding state notation is
\[
  \ket{\mathcal O(x)}
  =
  \ee^{-\ii x^\mu\widehat P_\mu}\ket{\mathcal O}.
\]

\section{Spin Sources and Contact Terms}
\label{sec:primary-spin-sources-contact-terms}

The primary module data must agree with the source-defined Ward identities.
Introduce a source \(\eta^a(x)\) valued in the dual representation
\(V_\rho^\vee\) and couple it by
\[
  \int\dd^D x\,\eta^a(x)\mathcal O_a(x).
\]
For a primary of dimension \(\Delta\), the flat-space source convention
associated with a conformal Killing vector \(\epsilon^\mu\) is
\begin{equation}
  \delta_\epsilon^{\rm src}\eta^a
  =
  \epsilon^\mu\partial_\mu\eta^a
  +
  (\Delta-D)\sigma_\epsilon\eta^a
  +
  \frac{\ii}{2}\omega_\epsilon^{\mu\nu}
  \eta^b(S_{\mu\nu})_b{}^a ,
  \label{eq:primary-source-transform}
\end{equation}
where
\[
  \sigma_\epsilon=\frac1D\partial_\mu\epsilon^\mu,
  \qquad
  \omega_{\epsilon,\mu\nu}
  =
  \frac12(\partial_\nu\epsilon_\mu-\partial_\mu\epsilon_\nu).
\]
The last term is the dual spin action.  The sign convention in
\eqref{eq:primary-source-transform} is the source convention of
Section~\ref{sec:cft-source-functionals-ward-contacts}; operator insertions
then acquire the contact operator
\begin{equation}
  \mathcal D_\epsilon\mathcal O_a(x)
  =
  -\epsilon^\mu(x)\partial_\mu\mathcal O_a(x)
  -\Delta\sigma_\epsilon(x)\mathcal O_a(x)
  -\frac{\ii}{2}\omega_\epsilon^{\mu\nu}(x)
  (S_{\mu\nu})_a{}^b\mathcal O_b(x).
  \label{eq:primary-contact-operator}
\end{equation}
Indeed, differentiating the source Ward identity with respect to the sources
and then setting the sources to zero gives, for
\[
  X=\mathcal O_{1,a_1}(x_1)\cdots\mathcal O_{n,a_n}(x_n),
  \qquad x_i\ne x_j ,
\]
the local current Ward identity
\begin{equation}
  \partial_\mu
  \left\langle J^\mu_\epsilon(x)X\right\rangle
  =
  \sum_{i=1}^n
  \delta^{(D)}(x-x_i)
  \left\langle
    \mathcal O_{1,a_1}(x_1)\cdots
    \mathcal D_\epsilon^{(i)}\mathcal O_{i,a_i}(x_i)\cdots
    \mathcal O_{n,a_n}(x_n)
  \right\rangle
  +
  \sigma_\epsilon(x)\mathcal A(x;X).
  \label{eq:primary-spin-current-ward}
\end{equation}
Here \(\mathcal D_\epsilon^{(i)}\) uses the spin representation and scaling
dimension of the \(i\)th insertion.  Equation
\eqref{eq:primary-spin-current-ward} is the spinful version of
\eqref{eq:local-conformal-current-ward}.  With the small-sphere orientation
used in \eqref{eq:variation-by-small-charge},
\begin{equation}
  \ii\,Q_\epsilon(S_x)\mathcal O_a(x)
  =
  \mathcal D_\epsilon\mathcal O_a(x).
  \label{eq:small-sphere-primary-action}
\end{equation}

\section{Infinitesimal Action at a General Point}

Let
\[
  \epsilon^\mu(x)
  =
  a^\mu+\omega^\mu{}_\nu x^\nu+\lambda x^\mu
  +c^\mu x^2-2(c\cdot x)x^\mu
\]
be a conformal Killing vector in \(D\ge3\).  Define
\[
  \sigma_\epsilon(x)=\frac1D\partial_\mu\epsilon^\mu(x),
  \qquad
  \omega_{\epsilon,\mu\nu}(x)
  =
  \frac12(\partial_\nu\epsilon_\mu-\partial_\mu\epsilon_\nu).
\]
For a primary operator, the infinitesimal variation generated by the charge is
the contact operator \(\mathcal D_\epsilon\):
\begin{equation}
  \delta_\epsilon\mathcal O_a(x)
  =
  -\epsilon^\mu(x)\partial_\mu\mathcal O_a(x)
  -\Delta\sigma_\epsilon(x)\mathcal O_a(x)
  -\frac{\ii}{2}\omega_{\epsilon}^{\mu\nu}(x)
  (S_{\mu\nu})_a{}^b\mathcal O_b(x).
  \label{eq:primary-infinitesimal-transform}
\end{equation}
The first term moves the insertion point.  The second term is the local scale
weight.  The third term rotates the spin indices by the orthogonal part of the
Jacobian.

For the four basic generators this gives
\begin{align}
  \delta_a\mathcal O(x)&=-a^\mu\partial_\mu\mathcal O(x),\\
  \delta_\lambda\mathcal O(x)&=-\lambda(x^\mu\partial_\mu+\Delta)\mathcal O(x),
  \label{eq:primary-translation-dilatation-transform}
\end{align}
and
\begin{align}
  \delta_\omega\mathcal O_a(x)
  &=
  -\omega^\mu{}_\nu x^\nu\partial_\mu\mathcal O_a(x)
  -\frac{\ii}{2}\omega^{\mu\nu}(S_{\mu\nu})_a{}^b\mathcal O_b(x),
  \\
  \delta_c\mathcal O_a(x)
  &=
  -\bigl(c^\mu x^2-2(c\cdot x)x^\mu\bigr)
  \partial_\mu\mathcal O_a(x)
  +2\Delta(c\cdot x)\mathcal O_a(x)
  \notag\\
  &\hspace{2.5em}
  -2\ii c^\mu x^\nu(S_{\mu\nu})_a{}^b\mathcal O_b(x).
  \label{eq:primary-sct-transform}
\end{align}
These formulae follow from \eqref{eq:primary-state-data},
\eqref{eq:operator-from-origin-translation}, and the conformal algebra.
In particular, the special-conformal formula is obtained by translating
\(\widehat K_\mu\ket{\mathcal O_a}=0\) from the origin to \(x\), or
equivalently by repeatedly commuting \(\widehat K_\mu\) through
\(\ee^{-\ii x^\rho\widehat P_\rho}\) in the Hermitian translation convention.
This derivation fixes the orbital, spin, and scaling terms in
\eqref{eq:primary-sct-transform} simultaneously.

\section{Finite Transformation Law}

Let
\[
  f:U\longrightarrow f(U)
\]
be an orientation-preserving conformal map between open subsets of
\(\R^D\).  Its Jacobian is
\[
  M_f(x)^\mu{}_\nu=\frac{\partial f^\mu}{\partial x^\nu}.
\]
The conformal condition says that there is a positive function
\(\Omega_f(x)\) such that
\begin{equation}
  M_f(x)^T M_f(x)=\Omega_f(x)^2\mathbf 1 .
  \label{eq:jacobian-scale-rotation-factorization}
\end{equation}
The orthogonal matrix
\[
  R_f(x)=\Omega_f(x)^{-1}M_f(x)
\]
is the local rotation.  Since \(f\) is orientation preserving,
\(\Omega_f(x)^D=\det M_f(x)\).  The factors obey the cocycle identities
\[
  \Omega_{f\circ h}(x)=\Omega_f(h(x))\Omega_h(x),
  \qquad
  R_{f\circ h}(x)=R_f(h(x))R_h(x).
\]
The finite primary transformation law is the integrated form of
\eqref{eq:primary-infinitesimal-transform} on the connected conformal group:
\begin{equation}
  \mathcal O'_a(f(x))
  =
  \Omega_f(x)^{-\Delta}
  \rho(R_f(x))_a{}^b\mathcal O_b(x).
  \label{eq:finite-primary-transform}
\end{equation}
For tensor representations this formula is the tensor transformation law with
the Jacobian decomposed into the scale factor \(\Omega_f(x)\) and the
orthogonal rotation \(R_f(x)\).  At coincident points the corresponding
statement is made at the level of the transformed source functional; the
displayed formula is its restriction to noncoincident insertions.

For a covariant rank-\(s\) primary tensor
\(\mathcal O_{\mu_1\cdots\mu_s}\), the same formula can be written without
abstract representation notation as
\begin{equation}
  \mathcal O'_{\mu_1\cdots\mu_s}(f(x))
  =
  \left(\det\frac{\partial f}{\partial x}\right)^{-(\Delta-s)/D}
  \frac{\partial x^{\nu_1}}{\partial f^{\mu_1}}\cdots
  \frac{\partial x^{\nu_s}}{\partial f^{\mu_s}}\,
  \mathcal O_{\nu_1\cdots\nu_s}(x).
  \label{eq:finite-primary-covariant-tensor}
\end{equation}
For a contravariant rank-\(s\) primary tensor,
\begin{equation}
  \mathcal O'^{\mu_1\cdots\mu_s}(f(x))
  =
  \left(\det\frac{\partial f}{\partial x}\right)^{-(\Delta+s)/D}
  \frac{\partial f^{\mu_1}}{\partial x^{\nu_1}}\cdots
  \frac{\partial f^{\mu_s}}{\partial x^{\nu_s}}\,
  \mathcal O^{\nu_1\cdots\nu_s}(x).
  \label{eq:finite-primary-contravariant-tensor}
\end{equation}
The different determinant powers in
\eqref{eq:finite-primary-covariant-tensor} and
\eqref{eq:finite-primary-contravariant-tensor} follow from
\(M_f=\Omega_f R_f\).  Each covariant index contributes one factor
\(\Omega_f^{-1}\) from \(M_f^{-1}\), each contravariant index contributes one
factor \(\Omega_f\) from \(M_f\), and the remaining determinant factor supplies
the primary weight \(\Omega_f^{-\Delta}\).

If the vacuum is invariant under the conformal transformation, correlation
functions of primaries obey
\begin{equation}
  \left\langle
    \prod_{i=1}^n\mathcal O_{i,a_i}(f(x_i))
  \right\rangle
  =
  \prod_{i=1}^n
  \Omega_f(x_i)^{-\Delta_i}
  \rho_i(R_f(x_i))_{a_i}{}^{b_i}
  \left\langle
    \prod_{i=1}^n\mathcal O_{i,b_i}(x_i)
  \right\rangle .
  \label{eq:primary-correlator-covariance}
\end{equation}
Equation \eqref{eq:primary-correlator-covariance} is the finite form of the
global Ward identities on the configuration space of distinct insertion
points, with contact extensions fixed by the same source convention as in
\eqref{eq:primary-spin-current-ward}.

\section{Inversion}

The inversion map
\[
  I(x)^\mu=\frac{x^\mu}{x^2}
\]
has
\[
  \Omega_I(x)=\frac1{x^2},
  \qquad
  R_I(x)^\mu{}_\nu
  =
  \delta^\mu{}_\nu-2\frac{x^\mu x_\nu}{x^2}.
\]
For a scalar primary,
\begin{equation}
  \mathcal O'(I(x))
  =
  |x|^{2\Delta}\mathcal O(x).
  \label{eq:scalar-primary-inversion}
\end{equation}
This formula, together with translations, rotations, and dilatations,
determines the separated-point covariance constraints used below.  The matrix
\(R_I(x)\) lies in \(O(D)\) and has determinant \(-1\).  For operators with
spin, inversion therefore requires a specified extension of the spin action
from \(SO(D)\) to the relevant disconnected orthogonal or pin action; once
that choice is made, \(R_I(x)\) acts on each spin index through that
representation.

\section{Descendants}

The descendants of a primary are obtained by applying derivatives.  Their
transformation laws are obtained by differentiating
\eqref{eq:finite-primary-transform} and by decomposing the result into
irreducible \(SO(D)\) representations.  The descendant transformation law
therefore contains no independent dynamical datum.  Once the primary
dimension, spin representation, and normalization are fixed, the action of the
whole conformal algebra on its multiplet is fixed.
